\section{Removing unused radix positions in nine-variable equations}
\label{sec:radix-origin}

This section gives three complete polynomial constructions in nine strictly
positive integers.  The improvement has two parts: an exact translation of
the finite coefficient layout, and a smaller evaluator for the translated
polynomial.  The translation is an existence-preserving change of the
encoding; the evaluator change is an identity of polynomials.  These two
claims should be kept distinct.

\begin{center}
\begin{tabular}{lrrr}
\hline
Construction & Variables & Exact degree & Binary operations\\
\hline
Original31, simple majorant &9&1,472,168&420\\
Original31, complete radical majorant &9&997,810&1,057\\
Compact54, simple majorant &9&1,626,248&428\\
\hline
\end{tabular}
\end{center}

The respective incoming comparisons are the saturated H1868 simple form
(1,474,688; 449 operations), its complete-radical form (999,518; 1,086), and
the compact54 saturated form (1,628,768; 434).  Thus the reductions relative
to those immediate predecessors are respectively 2,520 degrees and 29
operations, 1,708 degrees and 29 operations, and 2,520 degrees and six
operations.  The 420-operation form also has lower degree and fewer
operations than the 434-operation compact form, but it has a larger
coefficient payload.  These are improvements in stated measures within the
inspected repository.  They do not establish global minimality or reduce
the number of variables below nine.

The first two predecessors are in
\texttt{ConZFC\_Saturated\_Certificates\_2026-10-24.zip}; the third is in
\texttt{con-zfc-saturation-and-projection-2026-10-23.zip}.
The repository snapshot is
\texttt{7e533c33819ee1f703b69f2778a767cdc069940d}.
The included source ledger identifies these exact archives by SHA-256.

\subsection{A translation lemma for sparse coefficient layouts}

Let $\rho_1,\ldots,\rho_m$ be quadratic polynomials in $\nu$ variables,
and suppose fixed coefficient polynomials $K_0(T),K_1(T),K_2(T)$, distinct
nonnegative weights $w_i$, and distinct nonnegative target positions $q_j$
satisfy
\begin{equation}
 [T^{q_j}]\sum_{d=0}^{2}K_d(T)U_{\mathbf p}(T)^d
       =2\rho_j(\mathbf p),\qquad
 U_{\mathbf p}(T)=\sum_{i=1}^{\nu}p_iT^{w_i}.
 \label{eq:origin-extraction}
\end{equation}
The equality is a polynomial identity in all source coordinates, including
coefficients of monomials absent from the desired residual.  It is stronger
than a numerical identity at selected source tuples.

\begin{lemma}[Translation of the layout]
\label{lem:radix-origin}
Let $r,s$ be nonnegative integers such that $w_i\ge r$, $q_j\ge s$, and
every occupied exponent $e$ of $K_d$ satisfies $e+dr-s\ge0$.
Set
\[
 w_i'=w_i-r,\quad q_j'=q_j-s,\quad
 K_d'(T)=\sum_e [T^e]K_d(T)\,T^{e+dr-s}.
\]
Then these are ordinary polynomial data and satisfy the same extraction
identities with primes on $w,q,K,U$.
If $W=\max_iw_i$ and
$H=\max_d(\deg K_d+dW)$, their new values are
$W'=W-r$ and $H'=H-s$.
\end{lemma}

\begin{proof}
As a temporary identity in Laurent polynomials,
$U_{\mathbf p}'=T^{-r}U_{\mathbf p}$ and
$K_d'=T^{dr-s}K_d$.  Therefore
\[
 \sum_d K_d'(U_{\mathbf p}')^d
       =T^{-s}\sum_dK_dU_{\mathbf p}^d.
\]
Coefficient extraction proves the claim.  The hypotheses ensure that every
exponent in the actual delivered data is nonnegative, so the construction
contains no division by an unknown.  Finally
$\deg K_d'+dW'=\deg K_d+dW-s$, proving the height identity.
\end{proof}

Apply the lemma with $(r,s)=(1,2)$.  Both inherited layouts have minimum
coordinate weight one, minimum occupied linear-kernel exponent one, and
minimum occupied quadratic-kernel exponent zero.  Their constant kernels
begin above exponent two.  Consequently the transformation reads
\begin{equation}
 w_i'=w_i-1,\qquad q_j'=q_j-2,\qquad
 K_2'=K_2,\quad K_1'=K_1/T,\quad K_0'=K_0/T^2.
 \label{eq:origin-shift}
\end{equation}
The last two expressions denote removal of initial zero coefficients;
the JSON tables and the final circuits use nonnegative exponents only.
No source coordinate, source equation, or numerical coefficient changes.

There is a useful restricted optimality statement.  Among translations of
this form with $r,s\ge0$ and nonnegative resulting exponents, $r\le1$
because the minimum original weight is one, and $s\le2r$ because the
quadratic kernel has a nonzero constant coefficient.  Thus $(1,2)$ is
maximal in both parameters for these fixed layouts.  This says nothing
about changing the weights nonuniformly, changing the source residuals,
or using a different encoding architecture.

The translated coordinate weights, in source-variable order, are
\[
\begin{split}
 &(0,6,60,35,15,84,19,66,42,10,33,12,95,79,28,49,\\
 &\hspace{12mm}1,108,109,111,115,90,52,54,26,100,91,8,106,36,59).
\end{split}
\]
For the original31 source, the target positions in row order are
\[
\begin{split}
 &(1139,26,1435,1067,1834,1352,1812,612,265,1406,1401,\\
 &\hspace{12mm}887,335,1585,411,671,1137,770,1129,1016,1053,83).
\end{split}
\]
The source is the inherited positive system named \texttt{gapN86}, with
31 coordinates and 22 quadratic residuals.  Its old coordinate maps are
$P_{\rm old}=P+\eta+\zeta$ and $N_{\rm old}=N+1$; its recorded elementary
row operations are retained unchanged.  The compact source is the
separate \texttt{gapNq85} system.  Its tables and whole compact54 coefficient
tuple are kept separately; substituting individual coefficient ports
between the two systems is not justified.

\subsection{Why a coordinate may occupy position zero}

The choice $w_i>0$ was not needed by the positive-domain bit-mask argument.
For distinct weights in $\{0,\ldots,W\}$, put
\[
 \ell=\sum_iB^{w_i},\qquad N_0=B^{W+1},\qquad
 M_0=N_0-1-(b-1)\ell.
\]
Assume $b$ and $B$ are powers of two, with $b\mid B$ and $b\le B$.
The radix-$B$ digits of $M_0$ are $B-b$ at selected positions and $B-1$
elsewhere.  For $0\le g<N_0$, binary disjointness $g\mathbin{\&}M_0=0$
is therefore equivalent to
\begin{equation}
 g=\sum_i z_iB^{w_i},\qquad 0\le z_i<b.
 \label{eq:origin-mask}
\end{equation}
This digit argument includes position zero without any modification: the
allowed bits there are precisely its bottom $\log_2b$ bits.  The recovered
positive source coordinates are $p_i=z_i+1$.

It is essential to check the next arithmetic interface rather than merely
observe that the digit argument works.  The pinned Lean source
\texttt{Diophantine/Paper1982/NineRatio.lean}, retained in the packet,
states \texttt{Jones1982.Nine.ratio\_sound} and
\texttt{ratio\_complete} with the numerical hypotheses
\[
 R,N\ge8,\quad 0<b\le R,N,
\]
and with the appropriate positive carrier, margin, central-binomial
and power-of-two conditions.  Neither theorem assumes $b\mid g$ or
$b\mid R+1$.  Hence the possible nonzero units digit introduced by
\eqref{eq:origin-mask} does not violate an inherited premise.  This point
is specific to the positive nine-variable interface; the signed
Sun/Lucas construction has additional residue normalizations and is not
changed here.

The all-one source tuple is not a common zero: its fifth residual is one,
independently of the four coefficient ports.  Thus every genuine source
zero has at least one $p_i>1$, and its packed $g$ is positive even with a
weight zero.  For all sufficiently large binary $b$, its digits satisfy
\eqref{eq:origin-mask} and $g<bB^W$.  The cofinal completeness property
survives the translation.

\subsection{The fully recalibrated original31 prefix}

For this paragraph all layout data are the translated original31 data.
Write $K_d$ for their three kernels, and put
\[
 \kappa=2134939,\qquad B=2^\kappa b^2,\qquad
 T_*=\sum_{j=1}^{22}B^{q_j},\qquad U_* =\ell+g.
\]
Define
\begin{align}
 G_*&=(K_2(B)U_*+K_1(B))U_*+K_0(B),\notag\\
 N_0&=B^{116},& M_0&=N_0-1-(b-1)\ell,\notag\\
 S_*&=g+N_0(2G_*+T_*),&
 T_c&=M_0+N_0(B-2)T_*,\notag\\
 L&=2bB^{1983},&
 R&=L(S_*+T_c+1)-S_*-1,\notag\\
 N&=2^{16012032}B^{1053},&t&=bB^{115}.
 \label{eq:origin-prefix}
\end{align}
Every exponent here is fixed.  The kernel degrees are $(1406,1726,1636)$,
so $W=115$, $H=1866$, and the largest target is $1834=H-32$.
The leading kernel coefficients remain $-6,2,1$.

The inherited coefficient bound is $M=2^{2134927}$, and
$B=4096Mb^2$.  The exact artifact checker also verifies this bound on
all actual fixed coefficients.  For positive $b,g$, one has
$U_*\ge B^{115}$ and $K_2(B)>B^{1636}/2$.
Relative to $B^{1636}U_*^2$, the possible negative linear and constant
contributions are bounded by $2M/B^{25}$ and $2M/B^{460}$.
Thus $G_*>0$ unconditionally.  The low mask lies in $[0,N_0)$, so
$S_*,T_c\ge0$ and
\[
 R=(L-1)(S_*+T_c+1)+T_c\ge L-1\ge8.
\]
Also $b\le R,N$ and $N\ge8$.

On the soundness range $g<4t$, let $Z_*=B^{1867}$.
The same geometric-series estimate as in the inherited sparse-carry proof
now gives
\[
 U_*<6bB^{115},\qquad
 \frac{2G_*}{Z_*}<\frac{86Mb^2}{B-1}<\frac1{32}.
\]
The two target-mask bounds are unchanged because their headroom $H-Q=32$
is unchanged.  Together with $g<N_0$ and $M_0<N_0$, they give
\[
 S_*+T_c<N_0Z_*=B^{1983}<L.
\]
Each source convolution coefficient has magnitude at most
$993Mb^2<B/4$.  The retained signed-digit carry lemma therefore applies
with exactly the same strict inequalities.  It identifies
$S_*\mathbin{\&}T_c=0$ with a positive common source zero.

For completeness, the threshold must be recalculated.  Write
$b=2^k$, $B=2^v$, where $v=\kappa+2k$, and $L=2^u$.
There are 116 low radix digits, 31 selected low positions, and 22 high
mask digits.  Consequently
\[
 \operatorname{pop}(T_c)=138v-31k-22,\qquad
 u=1+k+1983v.
\]
The two-block population identity gives
\[
 \operatorname{pop}(R)=u+\operatorname{pop}(T_c)-\tau,
 \quad
 \tau=\operatorname{pop}(S_*)+\operatorname{pop}(T_c)
       -\operatorname{pop}(S_*+T_c)\ge0,
\]
with $\tau=0$ exactly for binary disjointness.  Now
\begin{align*}
 u+\operatorname{pop}(T_c)
   &=2121\kappa+4212k-21\\
   &=2\bigl(16012032+1053(\kappa+2k)\bigr)
     =2\log_2N,
\end{align*}
since $15\kappa-21=32024064$.
Legendre's formula therefore yields
\begin{equation}
 N^2\mid\binom{2R}{R}
 \quad\Longleftrightarrow\quad S_*\mathbin{\&}T_c=0.
 \label{eq:origin-calibration}
\end{equation}
The argument includes $b=1$.  At that value $M_0=N_0-1$ and
$0<g<N_0$, so disjointness fails.  There is no exceptional accepted
assignment at the endpoint.

For the compact54 variant replace the source tables and the whole native
coefficient tuple, and use
\[
 \kappa=864026,\quad H=2080,\quad Q=2048,\quad
 L=bB^{2197},\quad N=2^{6480184}B^{1160},\quad t=bB^{115}.
\]
The kernel degrees are $(1625,1940,1850)$, the dominance gaps are 25 and
455, and the coefficient bound is $2^{864014}$.  Its population calculation
is $u+\operatorname{pop}(T_c)=2335\kappa+4640k-22=2\log_2N$.
Thus the same proof applies.  The exact source audit checks all 22
identities separately for each of these two genuinely different sources.

\subsection{The actual nine-variable polynomials}

The unknowns are
\[
       b,g,h,s,w,C,i,j,n\in\mathbb Z_{>0}.
\]
All other names in this subsection are abbreviations, not quantified
variables.  Use either certified prefix above and define
\begin{align*}
 U&=N^2w,&Y&=N^2s,\\
 M&=RY,&P&=2M^2U,\\
 K&=R+1+h(P-1),&A&=M(U+1),\\
 B_1&=2R+1,&W_1&=bw,\\
 J&=4A-5,&Q&=3W_1C-2(W_1^2-1),\\
 m&=K^2(4t-g)-16t(C-KY)^2,&\Omega&=Jm,\\
 E_0&=J(n+W_1^2)+Q,&V_0&=PK,\\
 a_0&=(P^2-1)K^2+1,&D&=(A^2-1)C^2+1,\\
 r_0&=iJC^2,&F&=(A^2-1)Dr_0^2+1,\\
 \Lambda&=CD+1-A,&G&=A+F\Lambda,\\
 Z&=B_1+(j-1)C,&&\\
 H_1&=C+FZ,&I&=(G^2-1)H_1^2+1.
\end{align*}
For either simple construction set
\[
        a_1=DFI,\qquad V_1=ADr_0GH_1.
\]
For the complete-radical construction set
\[
 \sigma=1024C^5i,\qquad a_1=\sigma^2DFI,\qquad
 V_1=\Pi(A,C,i,Z),
\]
where $\Pi$ is the fully specified inherited 333-term integer polynomial
in \path{data/saturated_corrected_P.json}.
Its exact identity certificate is retained; it is not an unspecified
choice of a majorant.  Finally set
\begin{equation}
 \Theta=E_0-\Omega(V_0+V_1),\qquad
 \boxed{\bigl(\Theta^2-\Omega^2(a_0+a_1)\bigr)^2
                   -4\Omega^4a_0a_1=0.}
 \label{eq:origin-nine}
\end{equation}
The three literal division-free DAGs, all fixed numerical leaves, and
finite coefficient reconstruction recipes are included in the packet.

\begin{theorem}[Preservation of the consistency encoding]
The three positive-domain equations \eqref{eq:origin-nine} have a zero
if and only if their respective inherited positive source systems have
a common zero.  Composing with the retained native-source and
machine/refutation calibration gives the same ordinary theorem
formalizable in ZFC:
\[
 \mathrm{ZFC}\vdash
 \bigl(\nexists(b,g,h,s,w,C,i,j,n)\in\mathbb Z_{>0}^9:
       \mathcal F=0\bigr)\ \longleftrightarrow\ \operatorname{Con}(\mathrm{ZFC}).
\]
\end{theorem}

\begin{proof}
The preceding arguments establish all the numerical prefix hypotheses,
its soundness range, cofinal completeness, and the exact central-binomial
calibration.  The source zero set is unchanged by
Lemma~\ref{lem:radix-origin}.  The saturated positive ProductPell finalizer
is retained unchanged.

For clarity, its proof order matters.  Unconditional strict bounds
$V_0>\sqrt{a_0}$ and $V_1>\sqrt{a_1}$ first force $m>0$ from a norm zero.
They then give the carrier range and strict separation of the two
radicands.  A nonzero rational sum of their signed square roots forces
each root to be rational, hence integral.  This supplies the square
conditions and $J\mid Q$.  The saturated ProductPell theorem recovers
$C=\psi_A(2R+1)$, after which the original numerical ratio theorem
recovers the binary base and \eqref{eq:origin-calibration}.  No inference
about square roots is used prematurely to obtain the margin.
Conversely a source zero supplies fresh ratio and Pell witnesses,
and the inherited integral choice of $n>0$ makes one norm factor vanish.
The recorded soundness and saturation proofs prove these steps in full.

For the simple majorant the identity
\[
 (ADr_0)^2-DF=D(Dr_0^2-1)>0
\]
implies the required strict bound after multiplication by $G^2H_1^2$.
For the complete majorant, the exact polynomial
$\Pi^2-\sigma^2DFI$ has 36,309 strictly positive coefficients after
$A=x+3,C=y+1,i=u+1,Z=v+1$, including a positive constant.
The translated polynomial $\Pi$ itself also has positive coefficients.
Thus the same majorant premise holds uniformly.  These are retained
algebraic certificates, rechecked exactly on the present parameters.

All new steps are elementary finite polynomial identities, integer digit
arguments, or arithmetic statements with the same formalizable inherited
interfaces.  The packet does not supply a new end-to-end Lean/Rocq proof
or an exported object-language ZFC derivation; its calibration dependencies
are explicit rather than inferred from finite testing.
\end{proof}

\subsection{Exact degrees and exact arithmetic-circuit checks}

For the original31 prefix,
\[
 r=\deg R=7931,\quad \nu=\deg N=2106,\quad
 a=\deg A=r+4\nu+2=16357.
\]
For compact54 these values are $8787,2320,18069$.
The positive leading source term is $R\sim 2^{\lambda+1}bB^{2E-1}$,
where $(E,\lambda)=(1983,1)$ or $(2197,0)$.
Thus these bounds are attained on the actual fixed-coefficient instances.
Write $\delta_0=V_0^2-a_0$, $\delta_1=V_1^2-a_1$, and
\[
 B_s=E_0^2-2E_0\Omega(V_0+V_1)+\Omega^2(\delta_0+\delta_1).
\]
The exact cancellation-resolved identity is
\begin{equation}
 \mathcal F=B_s^2+4\Omega^2V_0V_1B_s
                  +4\Omega^4(\delta_0V_1^2+a_0\delta_1).
 \label{eq:origin-degree-identity}
\end{equation}
One has
\[
 \deg\Omega=5a+2W+3,\quad \deg V_0=3a+r+1,\quad
 \deg\delta_0=3a+r+2,\quad \deg E_0=a+4.
\]
For the simple majorant,
$\deg V_1=18a+r+24$ and $\deg\delta_1=34a+2r+48$.
The uniquely dominant term of \eqref{eq:origin-degree-identity} is
$\Omega^4\delta_1^2$, giving
\[
                  \deg\mathcal F=88a+4r+8W+108.
\]
For the complete majorant, exact expansion gives
\[
 \deg V_1=18a+r+30,\qquad \deg\delta_1=18a+2r+60.
\]
Their unique leading monomials before composition are
\[
1048576A^{18}C^{24}i^6Z,\qquad
624687054848A^{18}C^{48}i^{12}Z^2.
\]
The four contribution degrees in \eqref{eq:origin-degree-identity} are
$948768,997810,989850,752426$.  The second is uniquely maximal, and its
leading coefficient is positive, giving degree $997810$.
Merely propagating maximum degrees through the unresolved norm would
not establish these exact values.

The evaluator uses one addition chain for every required power of $B$.
The selected circuits evaluate $\ell$ and $T_*$ as sums of shared powers,
and evaluate the three coefficient kernels in Horner form.  This saves
26 multiplications for the original31 forms compared with using Horner
form for both masks; it saves 25 for the compact profile.  The common
power chain is a finite table of equalities $e=u+v$; its use costs one
multiplication per live chain step.  Fixed numeral construction is excluded
from this convention, but all live binary additions, subtractions,
squares, and multiplications by numerals are counted.

The supplied verifier does not import the builder.  It proves 54 or 55
exact polynomial identities at the named registers of each DAG, including
all complete-majorant coefficients.  It separately checks the variable
set, acyclicity, liveness, operation counts, fixed-leaf recipes, and the
final output.  A second independently implemented sparse ring checks all
22 source identities including absent monomials, all 20,466 original31
and 22,450 compact convolution interactions, the coefficient bounds, the
positive majorant certificate, the population exponents, and the resolved
degree.  Eight complete native coefficient byte strings are reconstructed
from the retained literal finite inputs.  These checks support exact
algebra and the numerical instances; they do not experimentally test
consistency or substitute for the unbounded semantic calibration.

The coefficient-size tradeoff is material.  The original31 tuple occupies
533,099 raw bytes in its four ports, versus 215,492 for compact54.
Its twelve live fixed leaves occupy 4,001,391 raw bytes and 53,739 gzip
bytes, versus eleven leaves occupying 1,565,018 raw bytes and 23,481 gzip
bytes for compact54.  The larger live-leaf total includes the fixed powers
of two and derived coefficient products.  These are distinct measures
from circuit length and from the size of the entire research archive.

\subsection{Further questions suggested by the construction}

The uniform translation is maximal for these fixed supports, but a joint
optimization of coordinate weights, target positions, row operations and
evaluator structure may improve all three circuits further.  Such a search
should allow weight zero from its beginning and score the final polynomial
rather than only the convolution height.  Another concrete target is a
Lean proof of the general layout-translation and zero-position mask lemmas,
composed with the existing ratio interfaces.  Finally, the exact register
certificates suggest a small checked circuit verifier that accepts the
whole finite DAG and the fixed numerical leaves without trusting the
arithmetic-circuit producer.  None of these questions entails a claim that
nine variables are necessary or that a different direct proof encoding
cannot give a simpler equation.
